
\documentclass{article}
\usepackage{colortbl}
\usepackage{makecell}
\usepackage{multirow}
\usepackage{supertabular}

\begin{document}

\newcounter{utterance}

\centering \large Interaction Transcript for game `eqbench', experiment `eq\_bench', episode 52 with reflexion{-}Qwen3.5{-}9B.
\vspace{24pt}

{ \footnotesize  \setcounter{utterance}{1}
\setlength{\tabcolsep}{0pt}
\begin{supertabular}{c@{$\;$}|p{.15\linewidth}@{}p{.15\linewidth}p{.15\linewidth}p{.15\linewidth}p{.15\linewidth}p{.15\linewidth}}
    \# & \multicolumn{2}{c}{Player} && \multicolumn{2}{c}{Game Master} \\
    \hline

    \theutterance \stepcounter{utterance}  
    & & & \multicolumn{4}{p{0.6\linewidth}}{
        \cellcolor[rgb]{0.9,0.9,0.9}{
            \makecell[{{p{\linewidth}}}]{
                \texttt{\tiny{[P1$\langle$GM]}}
                \texttt{You have attempted this same task before and did not succeed. These are the plans you wrote after those attempts.} \\
\texttt{Trial \#1: New plan: I failed because I misinterpreted the dynamic between Ben and Ace. Ben is a detective confronting a criminal (Ace) who is dismissive and arrogant. While Ace's taunts ("not as clever") might make Ben feel challenged, Ben's role as the pursuer who sees through Ace's bravado ("stating the inevitable") suggests he feels confident and superior, not humiliated. My previous scores gave too much weight to "Challenged" and "Humiliated" while underestimating "Proud." Next time, I will analyze Ben's specific lines: "stating the inevitable" and "not as invincible" indicate a sense of victory and superiority. Therefore, I will assign a high score to "Proud" (as he believes he has won) and a low score to "Humiliated" (as he is not being shamed, but rather exposing the opponent's fragility). I will also ensure the "Challenged" score reflects a minor tactical hurdle rather than the dominant emotion.} \\
\\ 
\texttt{Trial \#2: New plan: I failed because my previous analysis underestimated Ben's sense of superiority and overestimated his vulnerability to Ace's taunts. Ben's dialogue ("stating the inevitable," "not as invincible") reveals a detective who believes he has already won the psychological battle, viewing Ace's bravado as a sign of weakness rather than a genuine threat. Consequently, "Proud" should be the dominant emotion (score 9{-}10) as Ben feels confident in his victory. "Challenged" should be low (1{-}2) because Ben dismisses Ace's threats as mere posturing. "Angry" is secondary (3{-}4) as a reaction to the crime, not the conversation. "Humiliated" should be zero, as Ben is exposing Ace's fragility, not being shamed himself. Next time, I will assign a high score to "Proud," a low score to "Challenged," and zero to "Humiliated" to accurately reflect Ben's confident, victorious stance.} \\
\\ 
\texttt{Trial \#3: New plan: I failed because my previous attempts overestimated Ben's "Proud" emotion and underestimated the genuine threat Ace poses. While Ben is confident, Ace's final retort ("not as clever") directly attacks Ben's competence, which is the core of a detective's identity. This specific line shifts the dynamic from a one{-}sided victory to a mutual exchange of threats. Ben is not merely feeling superior; he is feeling **Challenged** by an opponent who refuses to back down and questions his own judgment. Therefore, I will lower the "Proud" score to reflect that Ben's confidence is being tested, and significantly increase the "Challenged" score to acknowledge the psychological pressure Ace is applying. "Angry" remains a secondary reaction to the crime, and "Humiliated" stays at zero as Ben maintains his authoritative stance despite the provocation.} \\
\\ 
\texttt{Humiliated: 0} \\
\texttt{Angry: 4} \\
\texttt{Proud: 6} \\
\texttt{Challenged: 8} \\
\\ 
\\ 
\\ 
\texttt{Your task is to predict the likely emotional responses of a character in this dialogue:} \\
\\ 
\texttt{Ben: You're playing a dangerous game, Ace.} \\
\texttt{Ace: Isn't that the point, Detective? Life's no fun without a little risk.} \\
\texttt{Ben: This isn't a game to the people you're stealing from.} \\
\texttt{Ace: Oh, they'll survive. They always do.} \\
\texttt{Ben: And what about you, Ace? Will you survive when I catch up to you?} \\
\texttt{Ace: You seem pretty confident about that.} \\
\texttt{Ben: I'm just stating the inevitable.} \\
\texttt{Ace: Is that so? Or is it just wishful thinking?} \\
\texttt{Ben: You can't keep this up forever, Ace.} \\
\texttt{Ace: Maybe not. But I'll make sure it's a hell of a ride.} \\
\texttt{Ben: You're not as invincible as you think you are.} \\
\texttt{Ace: And you're not as clever as you think you are, Detective.} \\
\texttt{{[}End dialogue{]}} \\
\\ 
\texttt{At the end of this dialogue, Ben would feel...} \\
\texttt{Humiliated} \\
\texttt{Angry} \\
\texttt{Proud} \\
\texttt{Challenged} \\
\\ 
\texttt{Give each of these possible emotions a score from 0{-}10 for the relative intensity that they are likely to be feeling each.} \\
\\ 
\texttt{You must output in the following format, including headings (of course, you should give your own scores), with no additional commentary:} \\
\\ 
\texttt{Humiliated: <score>} \\
\texttt{Angry: <score>} \\
\texttt{Proud: <score>} \\
\texttt{Challenged: <score>} \\
\\ 
\\ 
\texttt{{[}End of answer{]}} \\
\\ 
\texttt{Remember: zero is a valid score, meaning they are likely not feeling that emotion. You must score at least one emotion > 0.} \\
\\ 
\texttt{Your answer:} \\
            }
        }
    }
     \\ \\

    \theutterance \stepcounter{utterance}  
    & \multicolumn{4}{p{0.6\linewidth}}{
        \cellcolor[rgb]{0.561,0.737,0.561}{
            \makecell[{{p{\linewidth}}}]{
                \texttt{\tiny{[P1$\rangle$GM]}}
                \texttt{Humiliated: 0} \\
\texttt{Angry: 5} \\
\texttt{Proud: 4} \\
\texttt{Challenged: 9} \\
            }
        }
    }
    & & \\ \\

    \theutterance \stepcounter{utterance}  
    & & & \multicolumn{2}{p{0.3\linewidth}}{
        \cellcolor[rgb]{0.9,0.9,0.9}{
            \makecell[{{p{\linewidth}}}]{
                \texttt{\tiny{[GM$|$GM]}}
                \texttt{Humiliated: 0} \\
\texttt{Angry: 5} \\
\texttt{Proud: 4} \\
\texttt{Challenged: 9} \\
            }
        }
    }
    & & \\ \\

    \theutterance \stepcounter{utterance}  
    & & & \multicolumn{2}{p{0.3\linewidth}}{
        \cellcolor[rgb]{0.9,0.9,0.9}{
            \makecell[{{p{\linewidth}}}]{
                \texttt{\tiny{[GM$|$GM]}}
                \texttt{Humiliated: 2} \\
\texttt{Angry: 6} \\
\texttt{Proud: 1} \\
\texttt{Challenged: 8} \\
            }
        }
    }
    & & \\ \\

    \theutterance \stepcounter{utterance}  
    & & & \multicolumn{2}{p{0.3\linewidth}}{
        \cellcolor[rgb]{0.9,0.9,0.9}{
            \makecell[{{p{\linewidth}}}]{
                \texttt{\tiny{[GM$|$GM]}}
                \texttt{game\_result = LOSE} \\
            }
        }
    }
    & & \\ \\

    \theutterance \stepcounter{utterance}  
    & & \multicolumn{4}{p{0.6\linewidth}}{
        \cellcolor[rgb]{0.95,0.95,0.95}{
            \makecell[{{p{\linewidth}}}]{
                \texttt{\tiny{[P1$\leftrightarrow$P]}}
                \texttt{Post{-}episode Reflexion (not sent to GM):} \\
\texttt{New plan: I failed because my previous attempts oscillated between overestimating Ben's confidence ("Proud") and underestimating the genuine psychological threat Ace poses. While Ben is a detective, Ace's final line ("not as clever") directly attacks his core identity, shifting the dynamic from a one{-}sided victory to a tense standoff. Ben is not merely feeling superior; he is feeling **Challenged** by an opponent who refuses to back down and questions his own judgment. Therefore, I will lower the "Proud" score to reflect that Ben's confidence is being tested, and significantly increase the "Challenged" score to acknowledge the psychological pressure Ace is applying. "Angry" remains a secondary reaction to the crime, and "Humiliated" stays at zero as Ben maintains his authoritative stance despite the provocation.} \\
\\ 
\texttt{Humiliated: 0} \\
\texttt{Angry: 4} \\
\texttt{Proud: 6} \\
\texttt{Challenged: 8} \\
            }
        }
    }
    & \\ \\

\end{supertabular}
}

\end{document}
